% GENERATED FILE. Edit canonical lessons, then run npm run generate:books.
% canonical-source-hash: fnv1a64:c65d782c9716be10
% canonical-lessons: SA-C207-divyati, SA-C207-pravahati, SA-C207-asvadayati, SA-C207-upadisati, SA-C207-sammanayati

\chapter{Gambles, Flows, Tastes, Advises, Honours}
\label{ch:sa-gambles-flows-tastes-advises-honours}
\hlchaptermodality{207}

\begin{chapteropening}
\textbf{By the end of this chapter:} \emph{I can say gambles, flows, tastes, advises and honours.}

Complete the last lesson of chapter 207: I can say gambles, flows, tastes, advises and honours.
\end{chapteropening}

\section[\texorpdfstring{\sk{दीव्यति} (dīvyati)}{दीव्यति (dīvyati)}]{\sk{दीव्यति} (dīvyati) — gambles, plays}

\label{lesson:SA-C207-divyati}



\begin{quote}
Before the new one: say the Sanskrit for shuts, then the Sanskrit for opens.
\end{quote}

\subsection*{You'll want to know: \sk{दीव्यति}}
\textbf{\sk{दीव्यति}} — \emph{dīvyati} — \textquotedblleft{}gambles, plays\textquotedblright{}.

\begin{grammarlens}[title={What you've built}]
One more word for this chapter.
\end{grammarlens}

\subsection*{Your turn}
\begin{itemize}
\raggedright
  \item \emph{Say it:} \emph{dīvyati}
  \item \emph{Say it:} \emph{dīvyati}, once more
  \item \emph{Say it:} say \emph{dīvyati}
  \item \emph{Recall:} say the Sanskrit for shuts, then the Sanskrit for opens, then say \emph{dīvyati} again
\end{itemize}

\begin{tcolorbox}[breakable,colback=teal!4,colframe=teal!35!black,title={Before you move on}]
What does \sk{दीव्यति} mean? (\textquotedblleft{}Gambles, plays\textquotedblright{}.) Say it once more.
\end{tcolorbox}
\section[\texorpdfstring{\sk{प्रवहति} (pravahati)}{प्रवहति (pravahati)}]{\sk{प्रवहति} (pravahati) — flows}

\label{lesson:SA-C207-pravahati}



\begin{quote}
Before the new one: say the Sanskrit for opens, then the Sanskrit for gambles.
\end{quote}

\subsection*{You'll want to know: \sk{प्रवहति}}
\textbf{\sk{प्रवहति}} — \emph{pravahati} — \textquotedblleft{}flows\textquotedblright{}.

\begin{grammarlens}[title={What you've built}]
One more word for this chapter.
\end{grammarlens}

\subsection*{Your turn}
\begin{itemize}
\raggedright
  \item \emph{Say it:} \emph{pravahati}
  \item \emph{Say it:} \emph{pravahati}, once more
  \item \emph{Say it:} say \emph{pravahati}
  \item \emph{Recall:} say the Sanskrit for opens, then the Sanskrit for gambles, then say \emph{pravahati} again
\end{itemize}

\begin{tcolorbox}[breakable,colback=teal!4,colframe=teal!35!black,title={Before you move on}]
What does \sk{प्रवहति} mean? (\textquotedblleft{}Flows\textquotedblright{}.) Say it once more.
\end{tcolorbox}
\section[\texorpdfstring{\sk{आस्वादयति} (āsvādayati)}{आस्वादयति (āsvādayati)}]{\sk{आस्वादयति} (āsvādayati) — tastes}

\label{lesson:SA-C207-asvadayati}



\begin{quote}
Before the new one: say the Sanskrit for gambles, then the Sanskrit for flows.
\end{quote}

\subsection*{You'll want to know: \sk{आस्वादयति}}
\textbf{\sk{आस्वादयति}} — \emph{āsvādayati} — \textquotedblleft{}tastes\textquotedblright{}.

\begin{grammarlens}[title={What you've built}]
One more word for this chapter.
\end{grammarlens}

\subsection*{Your turn}
\begin{itemize}
\raggedright
  \item \emph{Say it:} \emph{āsvādayati}
  \item \emph{Say it:} \emph{āsvādayati}, once more
  \item \emph{Say it:} say \emph{āsvādayati}
  \item \emph{Recall:} say the Sanskrit for gambles, then the Sanskrit for flows, then say \emph{āsvādayati} again
\end{itemize}

\begin{tcolorbox}[breakable,colback=teal!4,colframe=teal!35!black,title={Before you move on}]
What does \sk{आस्वादयति} mean? (\textquotedblleft{}Tastes\textquotedblright{}.) Say it once more.
\end{tcolorbox}
\section[\texorpdfstring{\sk{उपदिशति} (upadiśati)}{उपदिशति (upadiśati)}]{\sk{उपदिशति} (upadiśati) — advises, instructs}

\label{lesson:SA-C207-upadisati}



\begin{quote}
Before the new one: say the Sanskrit for flows, then the Sanskrit for tastes.
\end{quote}

\subsection*{You'll want to know: \sk{उपदिशति}}
\textbf{\sk{उपदिशति}} — \emph{upadiśati} — \textquotedblleft{}advises, instructs\textquotedblright{}.

\begin{grammarlens}[title={What you've built}]
One more word for this chapter.
\end{grammarlens}

\subsection*{Your turn}
\begin{itemize}
\raggedright
  \item \emph{Say it:} \emph{upadiśati}
  \item \emph{Say it:} \emph{upadiśati}, once more
  \item \emph{Say it:} say \emph{upadiśati}
  \item \emph{Recall:} say the Sanskrit for flows, then the Sanskrit for tastes, then say \emph{upadiśati} again
\end{itemize}

\begin{tcolorbox}[breakable,colback=teal!4,colframe=teal!35!black,title={Before you move on}]
What does \sk{उपदिशति} mean? (\textquotedblleft{}Advises, instructs\textquotedblright{}.) Say it once more.
\end{tcolorbox}
\section[\texorpdfstring{\sk{सम्मानयति} (sammānayati)}{सम्मानयति (sammānayati)}]{\sk{सम्मानयति} (sammānayati) — honours}

\label{lesson:SA-C207-sammanayati}



\begin{quote}
Before the new one: say the Sanskrit for tastes, then the Sanskrit for advises.
\end{quote}

\subsection*{You'll want to know: \sk{सम्मानयति}}
\textbf{\sk{सम्मानयति}} — \emph{sammānayati} — \textquotedblleft{}honours\textquotedblright{}.

\begin{grammarlens}[title={What you've built}]
One more word for this chapter.
\end{grammarlens}

\subsection*{Your turn}
\begin{itemize}
\raggedright
  \item \emph{Say it:} \emph{sammānayati}
  \item \emph{Say it:} \emph{sammānayati}, once more
  \item \emph{Say it:} say \emph{sammānayati}
  \item \emph{Recall:} say the Sanskrit for tastes, then the Sanskrit for advises, then say \emph{sammānayati} again
\end{itemize}

\begin{tcolorbox}[breakable,colback=teal!4,colframe=teal!35!black,title={Before you move on}]
What does \sk{सम्मानयति} mean? (\textquotedblleft{}Honours\textquotedblright{}.) Say it once more.
\end{tcolorbox}
